\chapter{Operating, maintaining, and evaluating an agent service}
\section{From a successful run to a service}
A notebook demonstration shows that a particular path can work. A service must handle valid inputs, malformed inputs, changing data, dependency outages, concurrent requests, and software updates. It also needs a maintainer who can explain failures and decide when to change or roll back the system.

For the equipment service, the operating contract includes supported request types, maximum latency, data-access scope, approval requirements, and fallback behavior. It should state what the system does when inventory is unavailable or policy evidence conflicts. A silent fallback to an invented answer may improve apparent responsiveness while violating the task.

The deployable unit is the complete configuration: model, prompts, tools, schemas, retrieval corpus, graph, policy, and evaluation set. Updating one component can invalidate assumptions in another. A new policy document can change eligibility; a new schema can make a previously accepted model response invalid; a different model can alter tool-selection behavior.

\section{Cost accounting}
For requests with known token rates, a simplified inference cost is:

\begin{equation}
C_{\text{model}}=\sum_{i=1}^{m}\left(r_{\text{in}}u_i+r_{\text{out}}v_i\right),
\end{equation}

where \(m\) is the number of model requests, \(u_i\) and \(v_i\) are input and output token counts, and the rates are expressed per token. Actual billing can distinguish cached input, model variants, reasoning usage, tools, and other categories. Use the provider's current definitions and observed usage for a real budget; the equation is a teaching model, not a quoted price schedule.

Total service cost also includes retrieval, storage, orchestration, monitoring, and human review. A cheaper model request may cause more repair calls or more reviewer work. Compare cost per attempted task and cost per successful task. If failures are common, quoting only the cost of a successful trace hides wasted work.

\section{Worked example: cost per success}
Suppose an invented experiment makes 100 attempts. The total measured resource cost is 20 units and 80 attempts succeed. Cost per attempt is 0.2 units; cost per success is \(20/80=0.25\) units. A second design costs 18 units and succeeds on 60 attempts. Its cost per attempt is lower, but its cost per success is \(18/60=0.3\) units.

The first design is more economical per successful outcome under these definitions. That does not settle the entire decision: latency, error severity, capacity, and user experience may differ. If one design makes unauthorized bookings, a favorable average cost does not make it acceptable.

A hard request count bounds one source of work but is not a currency cap. Requests can have different input and output sizes. For an actual spending limit, combine request bounds, token limits, observed usage, and provider-side account controls where available. State which bounds are exact and which are estimates.

\section{Latency and capacity}
Latency is the time one user waits; throughput is the amount of work completed per unit time. Parallel branches can lower latency while increasing simultaneous demand. At high load, queues can dominate the service's response time even when individual tool calls are fast.

Measure distributions rather than only means. A small number of very slow attempts can matter greatly to users. Report a suitable percentile alongside the median and failure rate, and state how timeouts are represented. Excluding timed-out attempts from latency summaries can make a struggling service look fast.

For a sequential agent loop, every extra model round adds waiting time. Before optimizing implementation details, inspect whether the round is necessary. Can a deterministic rule replace a model call? Can independent reads run together? Can a response be streamed while a later nonessential task continues? Each optimization must preserve the task contract and produce an honest status.

\section{Caching and freshness}
Caching reuses a prior result. It is useful when the result remains valid for the new request and caller. A cache key should include the inputs and relevant versions that determine validity. For a policy answer, that may include the document version. For a personalized result, it may include authorization scope or user identity where appropriate.

Inventory is dynamic. A cached availability observation may help browsing but cannot guarantee availability at commit time. The reservation service must revalidate. Likewise, a cached answer derived from private evidence must not be served to an unauthorized user simply because the question text matches.

A cache needs an invalidation rule or time-to-live grounded in the data's meaning. “Cache everything for an hour” is not a universal optimization. Measure hit rate and stale-result failures, and retain the evidence needed to explain why a cached result was considered valid.

\section{Retries, circuit breakers, and reconciliation}
A retry repeats work after a failure. Retry only when the failure category and operation semantics justify it. A malformed request should be corrected first. An authorization denial should remain denied. A transient read failure may merit a small bounded retry with delay. Retrying an effect requires idempotency or a reconciliation procedure.

A circuit breaker temporarily stops calls to a dependency after a defined failure pattern, allowing the service to return a controlled unavailable status rather than repeatedly overload a failing component. Its thresholds and recovery behavior are operating choices that need testing. A fallback must preserve the meaning of the task: it may offer an unverified suggestion, but it must not label it as a confirmed reservation.

Reconciliation compares the application's recorded state with the authoritative domain state after uncertainty. If a commit timed out, query by operation ID before creating another booking. If the service cannot determine whether the effect occurred, escalate that uncertainty rather than convert it to an ordinary “not booked” answer.

\section{Release gates and rollback}
A release gate defines evidence required before exposing a change. Run regression cases, held-out tasks, negative tests, and provider integration checks appropriate to the changed component. Compare complete system versions rather than only a model name. A prompt update that fixes one case can change unrelated tool behavior.

A staged rollout limits exposure while gathering operational evidence. The application needs a way to select versions and return to a known configuration. Rollback of software does not undo external effects already committed; those need a separate operational response. Keep these two meanings of recovery distinct.

Assign an owner for the corpus, tool contracts, approval policy, and evaluation set. If nobody owns policy freshness, adding better retrieval cannot solve stale evidence. If nobody maintains tests after a schema change, a green test run may validate an obsolete interface.

\section{What recent research can and cannot tell us}
The exploratory scientific-computing field report describes projects where verification and stewardship matter alongside model assistance \cite{scientific-report}. OpenAI's observational research-acceleration post discusses activity measures and their interpretation \cite{acceleration-report}. These sources motivate careful outcome validation and maintenance ownership. They do not estimate a universal productivity gain for students or for the equipment service.

The frontier is therefore a moving research question, not a list of capabilities to memorize. A new model or framework may change what is feasible, but the course's comparison method remains useful: define the task, identify the system boundary, measure complete attempts, inspect failures, and state what the evidence supports. Keep a dated reading log so later readers can distinguish a historical finding from a current interface claim.

\section{The capstone as an operating argument}
The capstone should demonstrate a bounded equipment-style service or another approved domain with the same engineering requirements. It needs a simpler baseline, evidence-backed answers, explicit limits, a checked approval path where effects are proposed, a recovery test, and a reproducible evaluation bundle. Offline operation must remain possible for the core assessment so a student's grade does not depend on a paid provider account.

The written report should explain why each framework is present. LangGraph makes the workflow and state transitions explicit. Gemini supplies model behavior through a configured interface. The OpenAI Agents SDK supplies an agent execution abstraction in the relevant lab path. None is an argument for adding unnecessary autonomy. A project that demonstrates that a simpler workflow is sufficient can be an excellent result.

Conclude with an operating note: who maintains the system, what signals trigger investigation, what happens when a dependency fails, and how results can be reproduced. This turns a working demonstration into a reviewable engineering claim.

\section{Exercises}
\begin{enumerate}
\item Forty attempts cost 12 teaching units and produce 30 successes. Compute cost per attempt and per success. State one important factor these ratios omit.
\item Design a cache key and invalidation policy for a public loan-policy answer. Explain why the same policy is insufficient for inventory commitment.
\item Specify which failures may be retried for a read operation and an effectful operation. Include the uncertain-commit case.
\item Write a release gate for changing the model used in Lab 7. Which local and live checks are required, and what remains uncertain after they pass?
\item Submit the capstone package and operating note. A reviewer must be able to trace every reported result to a saved attempt and rerun the offline path.
\end{enumerate}

\section{Further study and laboratory connection}
Use \cite{scientific-report,acceleration-report} as evidence-quality readings. Follow the capstone and submission templates supplied with the course. The final assessment rewards functional correctness, negative-case and recovery tests, controlled evaluation, and reproducibility; it does not require a positive result for any vendor or architecture.

